% intro.tex

\section{Introduction}
\label{sec:intro}

Pedestrian crossing-intention prediction -- estimating, from a short
observation window of an ego-vehicle's onboard sensors, whether a
visible pedestrian is about to cross the vehicle's path -- is a
safety-critical component of autonomous-driving perception stacks, and
has accumulated a substantial architecture literature built around two
benchmark datasets, PIE~\cite{rasouli2019pie} and
JAAD~\cite{rasouli2017jaad}. Almost all of this literature, from the
original SF-GRU baseline~\cite{rasouli2020sfgru} through
PCPA~\cite{kotseruba2021pcpa} to recent transformer- and Mamba-based
architectures~\cite{adapt2026,trajfusionnet2026,efficientpie2025},
reports results trained and tested \emph{within} a single dataset.
Gesnouin~et~al.~\cite{gesnouin2022} were, to our knowledge, the first to
test the natural follow-up question -- does a model trained on one of
these datasets generalize to the other? -- and found that it largely
does not: cross-dataset AUC collapses to near the level of a model with
no signal at all. No architecture paper we are aware of published since
has re-tested this.

This paper is an extended diagnosis of that collapse, carried out on
the SF-GRU architecture, together with one partial, principled fix.
Three things motivate treating this as a diagnosis-first paper rather
than a methods paper. First, while investigating cross-dataset transfer
we discovered that JAAD's standard data loader silently mislabels the
majority of samples under its default configuration
(Section~\ref{sec:method-labelbug}): pedestrians lacking a full
behavioral annotation are assigned a hard-coded, definitionally
uninformative label rather than being excluded, unless the loader's
\texttt{sample\_type} argument is set explicitly to the literature's
\texttt{JAAD\_behavior} convention. This is orthogonal to, and
independent of, cross-dataset generalization -- it affects in-domain
JAAD results too -- but every cross-dataset result in this paper (and,
we suspect, some fraction of prior JAAD-based work using this loader)
depends on getting it right, so we treat identifying and correcting it
as a contribution in its own right, not merely a footnote to our main
experiments.

Second, having corrected the label issue and confirmed the
cross-dataset collapse persists, we ran a falsification campaign
against the most natural explanations before proposing a fix. PIE and
JAAD were filmed with different, largely uncalibrated camera setups; a
plausible hypothesis is that cross-dataset failure is substantially a
recoverable geometric mismatch. We test this directly with a real
metric-space ground-plane reprojection using PIE's published camera
calibration and researched estimates for JAAD, together with a 20-point
sensitivity sweep over JAAD's unknown camera parameters
(Section~\ref{sec:results-geometry}), and find no evidence for it:
geometric normalization does not recover performance, and no plausible
camera-parameter choice for JAAD changes this conclusion. We similarly
test, and rule out, video-level train/test leakage as an explanation
for a separate, related phenomenon -- SF-GRU's sensitivity to which
train/test partition of a dataset is used
(Section~\ref{sec:results-leakage}). Both are reported as negative
results because we think ruling out plausible-sounding explanations is
useful independent of whether we found the true cause.

Third, motivated by a modality-zeroing ablation localizing
cross-dataset sensitivity to the visual scene-context branch rather
than raw trajectory coordinates, we apply domain-adversarial training
(a gradient-reversal domain classifier~\cite{ganin2015dann}) combined
with a contrastive scale-invariance term~\cite{chen2020simclr} on the
box-trajectory representation. This is, to our knowledge, the first
application of domain-adversarial training to cross-dataset pedestrian
crossing-intention transfer. It produces a real, seed-confirmed
improvement over the uncorrected baseline in both transfer directions
(PIE~$\to$~JAAD AUC $0.498\to0.621$; JAAD~$\to$~PIE AUC
$0.402\to0.613$; $n=6$ seeds each at our best configuration) -- but we
report this improvement without overstating it: absolute cross-dataset
AUC remains far below typical in-domain performance for this model
family, and we consider the gap that remains after our best fix to be
at least as important a finding as the gap that our fix closes. We
therefore probe \emph{why} the fix works, using an independent
representation-separability probe distinct from the domain classifier
used during training. We find a dissociation, not a correlation,
between probe-measured domain separability and task performance
(Section~\ref{sec:results-probe}): the configurations that improve
cross-dataset AUC also make the fused representation \emph{more}, not
less, separable by an independent probe, while a control that reduces
separability (same-class mixup) also \emph{reduces} AUC. This suggests
domain-adversarial training's benefit here does not operate through the
literal representational invariance it is nominally designed to
produce. We further test this recipe on two further, structurally
distinct fusion architectures: box-anchored cross-attention
(Section~\ref{sec:method-crossattn}) and uncertainty-weighted fusion
(Section~\ref{sec:method-uncertainty}, chosen because it is the
strongest \emph{unadapted} baseline architecture we find). On
cross-attention, the AUC improvement holds in one direction and is not
distinguishable from baseline in the other; probing both reveals the
separability increase itself tracks which direction improves -- large
where GRL + contrastive helps, negligible where it does not -- rather
than occurring uniformly regardless of outcome
(Section~\ref{sec:results-crossattn}). On uncertainty-weighted fusion,
the recipe does not merely fail to help: it significantly
\emph{degrades} the architecture's strong unadapted baseline in one
direction and trends down, non-significantly, in the other
(Section~\ref{sec:results-uncertainty}). We treat the combination as
evidence for, not against, the generic-regularizer account -- a
literal representation-alignment mechanism gives no obvious reason for
its effect on separability to scale with whether it helps or hurts the
task, nor for the same objective to reliably help one architecture and
reliably hurt another -- while being explicit that our fix is a
reproducible, useful result for SF-GRU specifically and not a
general-purpose recipe: of three architectures tested, it helps in
both transfer directions on only one.

In summary, this paper's contributions are:

\begin{enumerate}
\item Identification and correction of a silent labeling artifact in
JAAD's default data-loading configuration, affecting the large
majority of samples in the default (\texttt{sample\_type=`all'})
setting, together with a re-audit of cross-dataset collapse under the
corrected labels confirming it is not an artifact of this bug.
\item A falsification campaign ruling out camera-geometry mismatch
(via real metric-space reprojection and a camera-parameter sensitivity
sweep) and video-level train/test leakage as explanations for,
respectively, cross-dataset collapse and split-sensitivity.
\item A modality-zeroing ablation localizing cross-dataset sensitivity
primarily to the visual scene-context modality rather than raw
trajectory coordinates.
\item A domain-adversarial and contrastive training recipe that
partially, reproducibly ($n=6$ seeds, both directions) mitigates
cross-dataset collapse on SF-GRU, together with an honest
characterization of how much of the original gap it leaves unclosed.
\item A mechanistic probe analysis showing that this recipe's benefit
is dissociated from -- not explained by -- the representational
invariance it is nominally trained to induce, using controls (mixup, a
trivial track-length probe, and two further fusion architectures) to
support this reading.
\item A test of the recipe on two additional, structurally distinct
fusion architectures showing it is not a general-purpose fix: it
partially replicates on one (helping in one transfer direction) and
significantly \emph{harms} the strongest unadapted baseline of the
other, evidence we read as further support for the generic-regularizer
account over a literal domain-invariance account.
\item A test of whether a large, generically-pretrained
vision-language backbone (LoRA-fine-tuned Qwen2-VL) does better at
cross-dataset transfer than our from-scratch results, including an
explicit checkpoint-selection robustness check given JAAD's small
validation split -- finding no clear advantage, a second, qualitatively
different negative result for the pretrained-backbone hypothesis.
\end{enumerate}

We view this paper's primary value as diagnostic: a small,
well-controlled model family on two well-studied benchmarks, examined
closely enough to find a real data bug, rule out two plausible-sounding
explanations, and show that a working partial fix does not work for
the reason it was designed to. We believe this kind of audit is
currently under-supplied relative to the rate at which new
architectures are proposed for these same two benchmarks
(Section~\ref{sec:related}), none of which re-test cross-dataset
transfer.
